\section{Conclusion}
\label{sec:conclusion}

This project successfully developed a comprehensive drug information retrieval system, evolving from a raw data pipeline to a sophisticated hybrid search application.

The development began with establishing a robust data foundation by ingesting and characterizing datasets from OpenFDA and DrugBank. We defined a clear conceptual model centered on the "Drug" entity and implemented a reproducible ETL pipeline to clean and structure this diverse data.

Building on this foundation, we focused on core retrieval capabilities using Apache Solr. We implemented and evaluated multiple retrieval configurations, establishing a strong baseline using standard text analysis techniques. Our manual evaluation framework allowed us to quantify the performance of different field weighting strategies.

Subsequently, we significantly advanced the state of the art of our system by implementing a Hybrid Search architecture. Our extensive evaluation of dense vector models identified \textbf{PubMedBERT (NeuML)} as the superior embedding model for this domain. By combining this semantic signal with BM25 keyword matching using Reciprocal Rank Fusion (RRF), we achieved a dramatic performance improvement, with the optimal configuration ($\alpha=0.5$) doubling the Mean Average Precision (MAP) compared to vector-only search.

Finally, we wrapped these backend advances in a modern, user-centric frontend. Features like \textbf{Smart Filtering}, \textbf{Result Clustering}, \textbf{Semantic Snippets}, and \textbf{Visual Chemical Structures} (via PubChem) directly address real-world user needs for precision and explainability. The system now stands as a complete, end-to-end solution for pharmaceutical information retrieval, offering both high-accuracy search and an intuitive, feature-rich discovery experience. Future work could explore personalized ranking based on user history and expanding the knowledge graph with drug-interaction data.
